
\documentclass[preprint,12pt]{elsarticle}
\usepackage{booktabs}
\usepackage{graphicx}
\usepackage{amsmath}
\begin{document}
\journal{Robotics and Autonomous Systems}
\begin{frontmatter}
\title{Robotics and Autonomous Systems: Official Elsevier Submission Preview}
\tnotetext[t1]{Imported official template baseline: elsarticle-template-num.tex.}
\author{A. Placeholder}
\author{B. Example}
\address{Department of Verified Templates, Placeholder Institute}
\begin{abstract}
This submission-style preview is rendered with the imported official Elsevier elsarticle package and arranged as a realistic author-manuscript. The page is intentionally generic in scientific content but explicit about frontmatter, section flow, figure and table treatment, and the guide-aligned template baseline selected for this journal.
\end{abstract}
\begin{keyword}
journal template \sep official Elsevier package \sep author manuscript \sep formatting verification
\end{keyword}
\end{frontmatter}
\section{Introduction}
This verified preview follows the official Elsevier package imported from the publisher LaTeX instructions page rather than relying only on the local TeX installation. The section layout is expanded so the author-manuscript rhythm is easier to inspect.
\section{Related Work}
The related-work section gives the preview a more realistic research-paper hierarchy and helps expose how Elsevier frontmatter transitions into the technical body.
\section{Methodology}
The methodology section introduces representative technical content and a single-column figure so the imported template can be checked under a realistic load.
\begin{equation}
J = \sum_{k=1}^{N} \left( \alpha \|e_k\|_2^2 + \beta \|u_k\|_2^2 \right).
\end{equation}
\begin{figure}[t]
\centering
\fbox{\parbox[c][2.1in][c]{0.84\linewidth}{\centering Elsevier Figure Placeholder\\[0.4em]\normalsize System architecture, workflow, or experimental scene}}
\caption{Single-column figure used to inspect the imported Elsevier author-manuscript layout, caption spacing, and float balance.}
\end{figure}
\section{Experiments}
The experiments section provides enough content to inspect table density, caption placement, and the relationship between quantitative discussion and the surrounding body text.
\subsection{Experimental Setup}
This paragraph acts as a placeholder for datasets, hardware conditions, baseline settings, and protocol details that would normally appear in an Elsevier engineering paper.
\subsection{Quantitative Results}
\begin{table}[t]
\centering
\caption{Quantitative comparison block used to inspect the imported Elsevier template under the journal-specific citation mode and frontmatter style.}
\begin{tabular}{lcc}
\toprule
Method & Score & Time \\
\midrule
Placeholder A & 0.91 & 1.4 \\
Placeholder B & 0.88 & 1.7 \\
Placeholder C & 0.84 & 2.0 \\
\bottomrule
\end{tabular}
\end{table}
\section{Conclusion and Discussion}
This page verifies the imported official Elsevier package path and displays the selected sample-manuscript variant: elsarticle-template-num.tex. The surrounding guide-for-authors rules still define the final journal-specific limits, declarations, artwork handling, and reference behavior at submission.
\section*{Template Display}
Imported package source: Elsevier LaTeX instructions. Selected manuscript baseline: elsarticle-template-num.tex. Guide-aligned citation mode for this preview: numeric.
\begin{thebibliography}{00}
\bibitem{ref1} A. Smith, B. Lee, 2025. Template-aware Elsevier manuscript preparation. Journal of Placeholder Studies 10, 1--9.
\bibitem{ref2} C. Garcia, D. Patel, 2026. Figure and table balance in author manuscripts. Example Engineering Journal 14, 22--31.
\end{thebibliography}
\end{document}
